
\documentclass{article}
\usepackage{colortbl}
\usepackage{makecell}
\usepackage{multirow}
\usepackage{supertabular}

\begin{document}

\newcounter{utterance}

\centering \large Interaction Transcript for game `matchit\_ascii', experiment `similar\_grid\_2', episode 2 with deepseek{-}v4{-}pro.
\vspace{24pt}

{ \footnotesize  \setcounter{utterance}{1}
\setlength{\tabcolsep}{0pt}
\begin{supertabular}{c@{$\;$}|p{.15\linewidth}@{}p{.15\linewidth}p{.15\linewidth}p{.15\linewidth}p{.15\linewidth}p{.15\linewidth}}
    \# & $\;$A & \multicolumn{4}{c}{Game Master} & $\;\:$B\\
    \hline

    \theutterance \stepcounter{utterance}  
    & & \multicolumn{4}{p{0.6\linewidth}}{
        \cellcolor[rgb]{0.9,0.9,0.9}{
            \makecell[{{p{\linewidth}}}]{
                \texttt{\tiny{[P1$\langle$GM]}}
                \texttt{أنت تشارك في لعبة تخمين تعاونية.} \\
\texttt{الهدف هو معرفة ما إذا كانت هذه الشبكة وصورة أخرى لا يستطيع رؤيتها إلا أنا، هي نفسها.} \\
\texttt{يرجى وصف شبكتك أولاً. ثم سأقدم وصفي ويمكننا أن نطرح على بعضنا أسئلة حول الشبكات لمعرفة ما إذا كانت هي نفسها. الآن ابدأ وصفك القصير للشبكة بـ "الوصف:" متبوعًا بالوصف. لا تضف أي شيء آخر. هذه هي شبكتك:} \\
\texttt{X ▢ ▢ ▢ X} \\
\texttt{X X ▢ X X} \\
\texttt{X ▢ X ▢ X} \\
\texttt{X ▢ ▢ ▢ X} \\
\texttt{X ▢ ▢ ▢ X} \\
            }
        }
    }
    & \\ \\

    \theutterance \stepcounter{utterance}  
    & \multicolumn{4}{p{0.6\linewidth}}{
        \cellcolor[rgb]{0.9,0.9,0.9}{
            \makecell[{{p{\linewidth}}}]{
                \texttt{\tiny{[P1$\rangle$GM]}}
                \texttt{الوصف: شبكة 5x5 تحتوي على رموز X و▢، حيث تشكل X نمطًا يشبه الحرف "H" أو "I" مع امتداد سفلي. الصف الأول: X في العمودين الأول والخامس، والبقية ▢. الصف الثاني: X في جميع الأعمدة ما عدا الثالث. الصف الثالث: X في الأعمدة الأول والثالث والخامس، مع ▢ في الثاني والرابع. الصفان الرابع والخامس: X في العمود الأول فقط، والباقي ▢.} \\
            }
        }
    }
    & & \\ \\

    \theutterance \stepcounter{utterance}  
    & & & \multicolumn{2}{p{0.3\linewidth}}{
        \cellcolor[rgb]{0.9,0.9,0.9}{
            \makecell[{{p{\linewidth}}}]{
                \texttt{\tiny{[GM$|$GM]}}
                \texttt{continue} \\
            }
        }
    }
    & & \\ \\

    \theutterance \stepcounter{utterance}  
    & & \multicolumn{4}{p{0.6\linewidth}}{
        \cellcolor[rgb]{0.9,0.9,0.9}{
            \makecell[{{p{\linewidth}}}]{
                \texttt{\tiny{[GM$\rangle$P2]}}
                \texttt{أنت تشارك في لعبة تخمين تعاونية.} \\
\texttt{الهدف هو معرفة ما إذا كانت هذه الشبكة وصورة أخرى لا يستطيع رؤيتها إلا أنا، هي نفسها.} \\
\texttt{يرجى وصف شبكتك أولاً. ثم سأقدم وصفي ويمكننا أن نطرح على بعضنا أسئلة حول الشبكات لمعرفة ما إذا كانت هي نفسها. الآن ابدأ وصفك القصير للشبكة بـ "الوصف:" متبوعًا بالوصف. لا تضف أي شيء آخر. هذه هي شبكتك:} \\
\texttt{X ▢ ▢ ▢ X} \\
\texttt{X ▢ ▢ X X} \\
\texttt{X X X ▢ X} \\
\texttt{X ▢ ▢ ▢ X} \\
\texttt{X ▢ ▢ ▢ X} \\
            }
        }
    }
    & \\ \\

    \theutterance \stepcounter{utterance}  
    & & & \multicolumn{4}{p{0.6\linewidth}}{
        \cellcolor[rgb]{0.9,0.9,0.9}{
            \makecell[{{p{\linewidth}}}]{
                \texttt{\tiny{[GM$\langle$P2]}}
                \texttt{الوصف: شبكة مربعة 5×5 تحتوي على رموز X و ▢. الصف الأول: X ثم ثلاثة ▢ ثم X. الصف الثاني: X ثم ▢ ثم ▢ ثم X ثم X. الصف الثالث: X ثم X ثم X ثم ▢ ثم X. الصف الرابع: X ثم ثلاثة ▢ ثم X. الصف الخامس: X ثم ثلاثة ▢ ثم X.} \\
            }
        }
    }
     \\ \\

    \theutterance \stepcounter{utterance}  
    & & & \multicolumn{2}{p{0.3\linewidth}}{
        \cellcolor[rgb]{0.9,0.9,0.9}{
            \makecell[{{p{\linewidth}}}]{
                \texttt{\tiny{[GM$|$GM]}}
                \texttt{continue} \\
            }
        }
    }
    & & \\ \\

    \theutterance \stepcounter{utterance}  
    & & \multicolumn{4}{p{0.6\linewidth}}{
        \cellcolor[rgb]{0.9,0.9,0.9}{
            \makecell[{{p{\linewidth}}}]{
                \texttt{\tiny{[GM$\rangle$P2]}}
                \texttt{هذا هوالوصف: شبكة 5x5 تحتوي على رموز X و▢، حيث تشكل X نمطًا يشبه الحرف "H" أو "I" مع امتداد سفلي. الصف الأول: X في العمودين الأول والخامس، والبقية ▢. الصف الثاني: X في جميع الأعمدة ما عدا الثالث. الصف الثالث: X في الأعمدة الأول والثالث والخامس، مع ▢ في الثاني والرابع. الصفان الرابع والخامس: X في العمود الأول فقط، والباقي ▢.} \\
\texttt{اطرح الآن سؤالًا لمعرفة جوانب جديدة من شبكتي قد تكون مختلفة عن شبكتك. ابدأ بـ "السؤال:" ولا تضف أي شيء آخر.} \\
            }
        }
    }
    & \\ \\

    \theutterance \stepcounter{utterance}  
    & & & \multicolumn{4}{p{0.6\linewidth}}{
        \cellcolor[rgb]{0.9,0.9,0.9}{
            \makecell[{{p{\linewidth}}}]{
                \texttt{\tiny{[GM$\langle$P2]}}
                \texttt{السؤال: في الصف الثالث من شبكتك، هل يحتوي على ثلاثة رموز X متتالية في الأعمدة الأول والثاني والثالث، ثم ▢ ثم X، أم أن توزيع X مختلف؟} \\
            }
        }
    }
     \\ \\

    \theutterance \stepcounter{utterance}  
    & & & \multicolumn{2}{p{0.3\linewidth}}{
        \cellcolor[rgb]{0.9,0.9,0.9}{
            \makecell[{{p{\linewidth}}}]{
                \texttt{\tiny{[GM$|$GM]}}
                \texttt{continue} \\
            }
        }
    }
    & & \\ \\

    \theutterance \stepcounter{utterance}  
    & & \multicolumn{4}{p{0.6\linewidth}}{
        \cellcolor[rgb]{0.9,0.9,0.9}{
            \makecell[{{p{\linewidth}}}]{
                \texttt{\tiny{[P1$\langle$GM]}}
                \texttt{هذا هوالوصف: شبكة مربعة 5×5 تحتوي على رموز X و ▢. الصف الأول: X ثم ثلاثة ▢ ثم X. الصف الثاني: X ثم ▢ ثم ▢ ثم X ثم X. الصف الثالث: X ثم X ثم X ثم ▢ ثم X. الصف الرابع: X ثم ثلاثة ▢ ثم X. الصف الخامس: X ثم ثلاثة ▢ ثم X.} \\
\texttt{السؤال: في الصف الثالث من شبكتك، هل يحتوي على ثلاثة رموز X متتالية في الأعمدة الأول والثاني والثالث، ثم ▢ ثم X، أم أن توزيع X مختلف؟ابدأ إجابتك بـ "الإجابة:" ولا تضف أي شيء آخر.} \\
            }
        }
    }
    & \\ \\

    \theutterance \stepcounter{utterance}  
    & \multicolumn{4}{p{0.6\linewidth}}{
        \cellcolor[rgb]{0.9,0.9,0.9}{
            \makecell[{{p{\linewidth}}}]{
                \texttt{\tiny{[P1$\rangle$GM]}}
                \texttt{الإجابة: لا، توزيع X في الصف الثالث من شبكتي مختلف. فهو: X في العمود الأول، ▢ في الثاني، X في الثالث، ▢ في الرابع، X في الخامس. أي أنه لا يحتوي على ثلاثة رموز X متتالية.} \\
            }
        }
    }
    & & \\ \\

    \theutterance \stepcounter{utterance}  
    & & & \multicolumn{2}{p{0.3\linewidth}}{
        \cellcolor[rgb]{0.9,0.9,0.9}{
            \makecell[{{p{\linewidth}}}]{
                \texttt{\tiny{[GM$|$GM]}}
                \texttt{continue} \\
            }
        }
    }
    & & \\ \\

    \theutterance \stepcounter{utterance}  
    & & \multicolumn{4}{p{0.6\linewidth}}{
        \cellcolor[rgb]{0.9,0.9,0.9}{
            \makecell[{{p{\linewidth}}}]{
                \texttt{\tiny{[P1$\langle$GM]}}
                \texttt{اطرح الآن سؤالًا لمعرفة جوانب جديدة من شبكتي قد تكون مختلفة عن شبكتك. ابدأ بـ "السؤال:" ولا تضف أي شيء آخر.} \\
            }
        }
    }
    & \\ \\

    \theutterance \stepcounter{utterance}  
    & \multicolumn{4}{p{0.6\linewidth}}{
        \cellcolor[rgb]{0.9,0.9,0.9}{
            \makecell[{{p{\linewidth}}}]{
                \texttt{\tiny{[P1$\rangle$GM]}}
                \texttt{السؤال: في شبكتك، ما هي الرموز الموجودة في الصف الثاني، وتحديداً في الأعمدة الثاني والثالث والرابع؟} \\
            }
        }
    }
    & & \\ \\

    \theutterance \stepcounter{utterance}  
    & & & \multicolumn{2}{p{0.3\linewidth}}{
        \cellcolor[rgb]{0.9,0.9,0.9}{
            \makecell[{{p{\linewidth}}}]{
                \texttt{\tiny{[GM$|$GM]}}
                \texttt{continue} \\
            }
        }
    }
    & & \\ \\

    \theutterance \stepcounter{utterance}  
    & & \multicolumn{4}{p{0.6\linewidth}}{
        \cellcolor[rgb]{0.9,0.9,0.9}{
            \makecell[{{p{\linewidth}}}]{
                \texttt{\tiny{[GM$\rangle$P2]}}
                \texttt{الإجابة: لا، توزيع X في الصف الثالث من شبكتي مختلف. فهو: X في العمود الأول، ▢ في الثاني، X في الثالث، ▢ في الرابع، X في الخامس. أي أنه لا يحتوي على ثلاثة رموز X متتالية.} \\
\texttt{السؤال: في شبكتك، ما هي الرموز الموجودة في الصف الثاني، وتحديداً في الأعمدة الثاني والثالث والرابع؟ابدأ إجابتك بـ "الإجابة:" ولا تضف أي شيء آخر.} \\
            }
        }
    }
    & \\ \\

    \theutterance \stepcounter{utterance}  
    & & & \multicolumn{4}{p{0.6\linewidth}}{
        \cellcolor[rgb]{0.9,0.9,0.9}{
            \makecell[{{p{\linewidth}}}]{
                \texttt{\tiny{[GM$\langle$P2]}}
                \texttt{الإجابة: في الصف الثاني من شبكتي، الأعمدة الثاني والثالث والرابع هي: ▢ ثم ▢ ثم X.} \\
            }
        }
    }
     \\ \\

    \theutterance \stepcounter{utterance}  
    & & & \multicolumn{2}{p{0.3\linewidth}}{
        \cellcolor[rgb]{0.9,0.9,0.9}{
            \makecell[{{p{\linewidth}}}]{
                \texttt{\tiny{[GM$|$GM]}}
                \texttt{continue} \\
            }
        }
    }
    & & \\ \\

    \theutterance \stepcounter{utterance}  
    & & \multicolumn{4}{p{0.6\linewidth}}{
        \cellcolor[rgb]{0.9,0.9,0.9}{
            \makecell[{{p{\linewidth}}}]{
                \texttt{\tiny{[GM$\rangle$P2]}}
                \texttt{اطرح الآن سؤالًا لمعرفة جوانب جديدة من شبكتي قد تكون مختلفة عن شبكتك. ابدأ بـ "السؤال:" ولا تضف أي شيء آخر.} \\
            }
        }
    }
    & \\ \\

    \theutterance \stepcounter{utterance}  
    & & & \multicolumn{4}{p{0.6\linewidth}}{
        \cellcolor[rgb]{0.9,0.9,0.9}{
            \makecell[{{p{\linewidth}}}]{
                \texttt{\tiny{[GM$\langle$P2]}}
                \texttt{السؤال: في شبكتك، هل الصفان الرابع والخامس متطابقان تماماً، حيث يحتوي كل منهما على X في العمود الأول فقط و▢ في باقي الأعمدة؟} \\
            }
        }
    }
     \\ \\

    \theutterance \stepcounter{utterance}  
    & & & \multicolumn{2}{p{0.3\linewidth}}{
        \cellcolor[rgb]{0.9,0.9,0.9}{
            \makecell[{{p{\linewidth}}}]{
                \texttt{\tiny{[GM$|$GM]}}
                \texttt{continue} \\
            }
        }
    }
    & & \\ \\

    \theutterance \stepcounter{utterance}  
    & & \multicolumn{4}{p{0.6\linewidth}}{
        \cellcolor[rgb]{0.9,0.9,0.9}{
            \makecell[{{p{\linewidth}}}]{
                \texttt{\tiny{[P1$\langle$GM]}}
                \texttt{الإجابة: في الصف الثاني من شبكتي، الأعمدة الثاني والثالث والرابع هي: ▢ ثم ▢ ثم X.} \\
\texttt{السؤال: في شبكتك، هل الصفان الرابع والخامس متطابقان تماماً، حيث يحتوي كل منهما على X في العمود الأول فقط و▢ في باقي الأعمدة؟ابدأ إجابتك بـ "الإجابة:" ولا تضف أي شيء آخر.} \\
            }
        }
    }
    & \\ \\

    \theutterance \stepcounter{utterance}  
    & \multicolumn{4}{p{0.6\linewidth}}{
        \cellcolor[rgb]{0.9,0.9,0.9}{
            \makecell[{{p{\linewidth}}}]{
                \texttt{\tiny{[P1$\rangle$GM]}}
                \texttt{الإجابة: نعم، الصفان الرابع والخامس في شبكتي متطابقان تماماً، حيث يحتوي كل منهما على X في العمود الأول فقط و▢ في باقي الأعمدة.} \\
            }
        }
    }
    & & \\ \\

    \theutterance \stepcounter{utterance}  
    & & & \multicolumn{2}{p{0.3\linewidth}}{
        \cellcolor[rgb]{0.9,0.9,0.9}{
            \makecell[{{p{\linewidth}}}]{
                \texttt{\tiny{[GM$|$GM]}}
                \texttt{continue} \\
            }
        }
    }
    & & \\ \\

    \theutterance \stepcounter{utterance}  
    & & \multicolumn{4}{p{0.6\linewidth}}{
        \cellcolor[rgb]{0.9,0.9,0.9}{
            \makecell[{{p{\linewidth}}}]{
                \texttt{\tiny{[P1$\langle$GM]}}
                \texttt{اطرح الآن سؤالًا لمعرفة جوانب جديدة من شبكتي قد تكون مختلفة عن شبكتك. ابدأ بـ "السؤال:" ولا تضف أي شيء آخر.} \\
            }
        }
    }
    & \\ \\

    \theutterance \stepcounter{utterance}  
    & \multicolumn{4}{p{0.6\linewidth}}{
        \cellcolor[rgb]{0.9,0.9,0.9}{
            \makecell[{{p{\linewidth}}}]{
                \texttt{\tiny{[P1$\rangle$GM]}}
                \texttt{السؤال: في شبكتك، ما هي الرموز الموجودة في الصف الثاني، العمودين الأول والثاني؟} \\
            }
        }
    }
    & & \\ \\

    \theutterance \stepcounter{utterance}  
    & & & \multicolumn{2}{p{0.3\linewidth}}{
        \cellcolor[rgb]{0.9,0.9,0.9}{
            \makecell[{{p{\linewidth}}}]{
                \texttt{\tiny{[GM$|$GM]}}
                \texttt{continue} \\
            }
        }
    }
    & & \\ \\

    \theutterance \stepcounter{utterance}  
    & & \multicolumn{4}{p{0.6\linewidth}}{
        \cellcolor[rgb]{0.9,0.9,0.9}{
            \makecell[{{p{\linewidth}}}]{
                \texttt{\tiny{[GM$\rangle$P2]}}
                \texttt{الإجابة: نعم، الصفان الرابع والخامس في شبكتي متطابقان تماماً، حيث يحتوي كل منهما على X في العمود الأول فقط و▢ في باقي الأعمدة.} \\
\texttt{السؤال: في شبكتك، ما هي الرموز الموجودة في الصف الثاني، العمودين الأول والثاني؟ابدأ إجابتك بـ "الإجابة:" ولا تضف أي شيء آخر.} \\
            }
        }
    }
    & \\ \\

    \theutterance \stepcounter{utterance}  
    & & & \multicolumn{4}{p{0.6\linewidth}}{
        \cellcolor[rgb]{0.9,0.9,0.9}{
            \makecell[{{p{\linewidth}}}]{
                \texttt{\tiny{[GM$\langle$P2]}}
                \texttt{الإجابة: في الصف الثاني، العمود الأول هو X والعمود الثاني هو ▢.} \\
            }
        }
    }
     \\ \\

    \theutterance \stepcounter{utterance}  
    & & & \multicolumn{2}{p{0.3\linewidth}}{
        \cellcolor[rgb]{0.9,0.9,0.9}{
            \makecell[{{p{\linewidth}}}]{
                \texttt{\tiny{[GM$|$GM]}}
                \texttt{continue} \\
            }
        }
    }
    & & \\ \\

    \theutterance \stepcounter{utterance}  
    & & \multicolumn{4}{p{0.6\linewidth}}{
        \cellcolor[rgb]{0.9,0.9,0.9}{
            \makecell[{{p{\linewidth}}}]{
                \texttt{\tiny{[GM$\rangle$P2]}}
                \texttt{اطرح الآن سؤالًا لمعرفة جوانب جديدة من شبكتي قد تكون مختلفة عن شبكتك. ابدأ بـ "السؤال:" ولا تضف أي شيء آخر.} \\
            }
        }
    }
    & \\ \\

    \theutterance \stepcounter{utterance}  
    & & & \multicolumn{4}{p{0.6\linewidth}}{
        \cellcolor[rgb]{0.9,0.9,0.9}{
            \makecell[{{p{\linewidth}}}]{
                \texttt{\tiny{[GM$\langle$P2]}}
                \texttt{السؤال: في شبكتك، هل الصف الأول يحتوي على X في العمود الأول، ثم ثلاثة ▢، ثم X في العمود الخامس، أي مطابق للصفين الرابع والخامس؟} \\
            }
        }
    }
     \\ \\

    \theutterance \stepcounter{utterance}  
    & & & \multicolumn{2}{p{0.3\linewidth}}{
        \cellcolor[rgb]{0.9,0.9,0.9}{
            \makecell[{{p{\linewidth}}}]{
                \texttt{\tiny{[GM$|$GM]}}
                \texttt{continue} \\
            }
        }
    }
    & & \\ \\

    \theutterance \stepcounter{utterance}  
    & & \multicolumn{4}{p{0.6\linewidth}}{
        \cellcolor[rgb]{0.9,0.9,0.9}{
            \makecell[{{p{\linewidth}}}]{
                \texttt{\tiny{[P1$\langle$GM]}}
                \texttt{الإجابة: في الصف الثاني، العمود الأول هو X والعمود الثاني هو ▢.} \\
\texttt{السؤال: في شبكتك، هل الصف الأول يحتوي على X في العمود الأول، ثم ثلاثة ▢، ثم X في العمود الخامس، أي مطابق للصفين الرابع والخامس؟ابدأ إجابتك بـ "الإجابة:" ولا تضف أي شيء آخر.} \\
            }
        }
    }
    & \\ \\

    \theutterance \stepcounter{utterance}  
    & \multicolumn{4}{p{0.6\linewidth}}{
        \cellcolor[rgb]{0.9,0.9,0.9}{
            \makecell[{{p{\linewidth}}}]{
                \texttt{\tiny{[P1$\rangle$GM]}}
                \texttt{الإجابة: نعم، الصف الأول في شبكتي مطابق للصفين الرابع والخامس، حيث يحتوي على X في العمود الأول، ثم ثلاثة ▢، ثم X في العمود الخامس.} \\
            }
        }
    }
    & & \\ \\

    \theutterance \stepcounter{utterance}  
    & & & \multicolumn{2}{p{0.3\linewidth}}{
        \cellcolor[rgb]{0.9,0.9,0.9}{
            \makecell[{{p{\linewidth}}}]{
                \texttt{\tiny{[GM$|$GM]}}
                \texttt{continue} \\
            }
        }
    }
    & & \\ \\

    \theutterance \stepcounter{utterance}  
    & & \multicolumn{4}{p{0.6\linewidth}}{
        \cellcolor[rgb]{0.9,0.9,0.9}{
            \makecell[{{p{\linewidth}}}]{
                \texttt{\tiny{[P1$\langle$GM]}}
                \texttt{اطرح الآن سؤالًا لمعرفة جوانب جديدة من شبكتي قد تكون مختلفة عن شبكتك. ابدأ بـ "السؤال:" ولا تضف أي شيء آخر.} \\
            }
        }
    }
    & \\ \\

    \theutterance \stepcounter{utterance}  
    & \multicolumn{4}{p{0.6\linewidth}}{
        \cellcolor[rgb]{0.9,0.9,0.9}{
            \makecell[{{p{\linewidth}}}]{
                \texttt{\tiny{[P1$\rangle$GM]}}
                \texttt{السؤال: في شبكتك، هل يوجد أي صف يحتوي على ثلاثة رموز X متجاورة أفقياً؟} \\
            }
        }
    }
    & & \\ \\

    \theutterance \stepcounter{utterance}  
    & & & \multicolumn{2}{p{0.3\linewidth}}{
        \cellcolor[rgb]{0.9,0.9,0.9}{
            \makecell[{{p{\linewidth}}}]{
                \texttt{\tiny{[GM$|$GM]}}
                \texttt{continue} \\
            }
        }
    }
    & & \\ \\

    \theutterance \stepcounter{utterance}  
    & & \multicolumn{4}{p{0.6\linewidth}}{
        \cellcolor[rgb]{0.9,0.9,0.9}{
            \makecell[{{p{\linewidth}}}]{
                \texttt{\tiny{[GM$\rangle$P2]}}
                \texttt{الإجابة: نعم، الصف الأول في شبكتي مطابق للصفين الرابع والخامس، حيث يحتوي على X في العمود الأول، ثم ثلاثة ▢، ثم X في العمود الخامس.} \\
\texttt{السؤال: في شبكتك، هل يوجد أي صف يحتوي على ثلاثة رموز X متجاورة أفقياً؟ابدأ إجابتك بـ "الإجابة:" ولا تضف أي شيء آخر.} \\
            }
        }
    }
    & \\ \\

    \theutterance \stepcounter{utterance}  
    & & & \multicolumn{4}{p{0.6\linewidth}}{
        \cellcolor[rgb]{0.9,0.9,0.9}{
            \makecell[{{p{\linewidth}}}]{
                \texttt{\tiny{[GM$\langle$P2]}}
                \texttt{الإجابة: نعم، الصف الثالث في شبكتي يحتوي على ثلاثة رموز X متجاورة أفقياً في الأعمدة الأول والثاني والثالث.} \\
            }
        }
    }
     \\ \\

    \theutterance \stepcounter{utterance}  
    & & & \multicolumn{2}{p{0.3\linewidth}}{
        \cellcolor[rgb]{0.9,0.9,0.9}{
            \makecell[{{p{\linewidth}}}]{
                \texttt{\tiny{[GM$|$GM]}}
                \texttt{continue} \\
            }
        }
    }
    & & \\ \\

    \theutterance \stepcounter{utterance}  
    & & \multicolumn{4}{p{0.6\linewidth}}{
        \cellcolor[rgb]{0.9,0.9,0.9}{
            \makecell[{{p{\linewidth}}}]{
                \texttt{\tiny{[GM$\rangle$P2]}}
                \texttt{الآن اتخذ قرارًا. ما رأيك: هل شبكتك والشبكة الأخرى الموصوفة متماثلتان؟  اكتب "القرار: نفس الشبكة" إذا كنت تعتقد أنهما نفس الشبكة أو "القرار: شبكات مختلفة" إذا كنت تعتقد أنهما شبكتان مختلفتان. لا تضف أي شيء آخر.} \\
            }
        }
    }
    & \\ \\

    \theutterance \stepcounter{utterance}  
    & & & \multicolumn{4}{p{0.6\linewidth}}{
        \cellcolor[rgb]{0.9,0.9,0.9}{
            \makecell[{{p{\linewidth}}}]{
                \texttt{\tiny{[GM$\langle$P2]}}
                \texttt{القرار: شبكات مختلفة} \\
            }
        }
    }
     \\ \\

    \theutterance \stepcounter{utterance}  
    & & & \multicolumn{2}{p{0.3\linewidth}}{
        \cellcolor[rgb]{0.9,0.9,0.9}{
            \makecell[{{p{\linewidth}}}]{
                \texttt{\tiny{[GM$|$GM]}}
                \texttt{continue} \\
            }
        }
    }
    & & \\ \\

    \theutterance \stepcounter{utterance}  
    & & & \multicolumn{2}{p{0.3\linewidth}}{
        \cellcolor[rgb]{0.9,0.9,0.9}{
            \makecell[{{p{\linewidth}}}]{
                \texttt{\tiny{[GM$|$GM]}}
                \texttt{success} \\
            }
        }
    }
    & & \\ \\

    \theutterance \stepcounter{utterance}  
    & & \multicolumn{4}{p{0.6\linewidth}}{
        \cellcolor[rgb]{0.9,0.9,0.9}{
            \makecell[{{p{\linewidth}}}]{
                \texttt{\tiny{[P1$\langle$GM]}}
                \texttt{الإجابة: نعم، الصف الثالث في شبكتي يحتوي على ثلاثة رموز X متجاورة أفقياً في الأعمدة الأول والثاني والثالث.} \\
\texttt{الآن اتخذ قرارًا. ما رأيك: هل شبكتك والشبكة الأخرى الموصوفة متماثلتان؟  اكتب "القرار: نفس الشبكة" إذا كنت تعتقد أنهما نفس الشبكة أو "القرار: شبكات مختلفة" إذا كنت تعتقد أنهما شبكتان مختلفتان. لا تضف أي شيء آخر.} \\
            }
        }
    }
    & \\ \\

    \theutterance \stepcounter{utterance}  
    & \multicolumn{4}{p{0.6\linewidth}}{
        \cellcolor[rgb]{0.9,0.9,0.9}{
            \makecell[{{p{\linewidth}}}]{
                \texttt{\tiny{[P1$\rangle$GM]}}
                \texttt{القرار: شبكات مختلفة} \\
            }
        }
    }
    & & \\ \\

    \theutterance \stepcounter{utterance}  
    & & & \multicolumn{2}{p{0.3\linewidth}}{
        \cellcolor[rgb]{0.9,0.9,0.9}{
            \makecell[{{p{\linewidth}}}]{
                \texttt{\tiny{[GM$|$GM]}}
                \texttt{continue} \\
            }
        }
    }
    & & \\ \\

    \theutterance \stepcounter{utterance}  
    & & & \multicolumn{2}{p{0.3\linewidth}}{
        \cellcolor[rgb]{0.9,0.9,0.9}{
            \makecell[{{p{\linewidth}}}]{
                \texttt{\tiny{[GM$|$GM]}}
                \texttt{success} \\
            }
        }
    }
    & & \\ \\

\end{supertabular}
}

\end{document}
